% Synthetisches Testdokument im Stil eines klassischen pdfLaTeX-Laborberichts
% (inputenc/T1, siunitx, booktabs, enumitem, tcolorbox, listings, thebibliography).
\documentclass[11pt,a4paper]{article}
\usepackage[utf8]{inputenc}
\usepackage[T1]{fontenc}
\usepackage[ngerman]{babel}
\usepackage[a4paper,margin=2.5cm]{geometry}
\usepackage{amsmath,amssymb}
\usepackage{graphicx}
\usepackage{booktabs}
\usepackage{float}
\usepackage{siunitx}
\usepackage{enumitem}
\usepackage[most]{tcolorbox}
\usepackage{listings}
\usepackage{xcolor}
\usepackage{hyperref}

\newcommand{\Rges}{R_{\mathrm{ges}}}
\newcommand{\messwert}[1]{\textbf{#1}}

\title{Laborbericht: Spannungsteiler}
\author{Erika Beispiel \and Max Muster}
\date{\today}

\begin{document}
\maketitle

\begin{abstract}
Dieser Bericht untersucht einen unbelasteten Spannungsteiler mit $\Rges = \SI{10}{\kilo\ohm}$.
\end{abstract}

\section{Aufgabenstellung}
\begin{itemize}[leftmargin=1.5em]
  \item Aufbau der Schaltung mit \SI{4.7}{\kilo\ohm} und \SI{5.3}{\kilo\ohm}.
  \item Messung bei \SI{12}{\volt} Versorgungsspannung.
\end{itemize}

\begin{enumerate}
  \item[a)] Berechnen Sie die Teilspannung.
  \item[b)] Vergleichen Sie mit dem \messwert{Messwert}.
\end{enumerate}

\vspace{0.5cm}

\section{Messung}
Die Ergebnisse zeigt Tabelle~\ref{tab:messung}.\quad Abweichungen sind gering.

\begin{table}[H]
\centering
\caption{Messwerte}
\label{tab:messung}
\small
\begin{tabular}{@{}lrr@{}}
\toprule
Größe & Rechnung & Messung \\
\midrule
$U_1$ & \SI{5.64}{\volt} & \SI{5.61}{\volt} \\
$U_2$ & \SI{6.36}{\volt} & \SI{6.40}{\volt} \\
\bottomrule
\end{tabular}
\end{table}

\begin{figure}[htbp]
  \centering
  \includegraphics[width=0.6\textwidth]{abbildungen/aufbau}
  \caption{Messaufbau mit $R_1$ und $R_2$}
  \label{fig:aufbau}
\end{figure}

\begin{tcolorbox}[title={Ergebnis}, colback=blue!5!white, colframe=blue!60!black]
\begin{align}
U_2 &= U \cdot \frac{R_2}{R_1 + R_2} \\
    &= \SI{6.36}{\volt}
\end{align}
\end{tcolorbox}

\section{Auswertung}
\begin{lstlisting}[language=Python, caption={Auswertung}]
r1, r2 = 4.7e3, 5.3e3
print(12 * r2 / (r1 + r2))
\end{lstlisting}

\begin{thebibliography}{9}
\bibitem{tietze} U.~Tietze, C.~Schenk: \textit{Halbleiter-Schaltungstechnik}. Springer.
\end{thebibliography}
\end{document}
